
\vspace{1cm}

\noindent  \textbf{Leia as instruções antes de fazer a lista:}

\vspace{0.3cm}
\noindent A resolução deve ser entregue em formato Jupyter Notebook ou Google Colab na plataforma Google Classroom, na tarefa designada e até o dia determinado.

\noindent  Não utilize nenhum módulo, método ou função já existente no Python exceto pela função \code{print}. Não é necessário testar se os dados passados por argumento são válidos a menos que pedido na questão.


\vspace{0.5cm}


\begin{enumerate}[label=\textbf{Q\arabic*.}]


  

    \item Escreva uma função em Python que recebe por argumento dois números inteiros e maiores que zero e calcula o somatório de todos os números múltiplos de três e/ou múltiplos de cinco dentro desta faixa (incluindo os números dos argumentos). Utilize a estrutura de repetição \code{for} para fazer a questão.

    \ifmostrarGabarito
\begin{minted}{python}
# Resposta:
def calcular_soma( valor_min, valor_max ):

    if (not type(valor_min)==int) or (not type(valor_max)==int):
        print("os valores passados por parametro precisam ser inteiros.")
        return -1
        
    # acerta o range de interação caso min > max (invertido)
    valor = valor_min 
    if valor_min > valor_max
        valor_max = valor_min
        valor_min = valor

    total = 0
    for valor in range( valor_min, valor_max+1, 1 ):
        if (valor % 3 == 0) or (valor % 5 == 0):
            total+=valor

    print(f"Somatório é de {valor}")
    return valor

\end{minted}
\fi


    \vspace{0.5cm}

    \item A função exponencial natural pode ser calculada utilizando a seguinte série:
    \[
    e^{x} = \sum_{n=0}^{m} \frac{x^{n}}{n!} = \frac{x^{0}}{0!} + \frac{x^{1}}{1!} + \frac{x^{2}}{2!} + \dots + \frac{x^{m}}{m!}
    \]
    Veja que $m$ representa o número de termos do somatório que são usados. Nessa questão você deve, necessariamente, fazer três funções em Python:
    \begin{itemize}
        \item uma para cálculo do fatorial (cujo nome deve terminar com \texttt{2a});
        \item uma para o cálculo da potência sem o uso do operador com dupla asterisco \texttt{**}, a potência deve ser calculada com \code{while} (cujo nome deve terminar com o número \texttt{2b});
        \item outra para cálculo da função exponencial (cujo nome deve terminar com o número \texttt{2c}).
    \end{itemize}
    A função para cálculo do fatorial (\texttt{2a}) recebe por argumento um número inteiro e maior ou igual a $0$ e retorna o fatorial desse número (lembrando que o fatorial de $0$ é $1$). A função para cálculo da potência (\texttt{2b}) recebe dois números inteiros e positivos por argumento e retorna o resultado do primeiro número elevado ao segundo número. A função para cálculo da exponencial (\texttt{2c}) recebe os números $x$ e $m$ por argumento (nessa ordem) e calcula a função exponencial de acordo com o somatório acima. A função de cálculo da função exponencial deve chamar a função fatorial para encontrar os denominadores e deve chamar a função de potência para cálculo dos numeradores. Considere que $x$ é um \code{float} e que $m$ é um inteiro positivo e maior ou igual a $0$. A função \texttt{2c} deve retornar o resultado da exponencial de acordo com o cálculo da série.

    \ifmostrarGabarito{
\begin{minted}{python}
# Resposta:
def calcular_fatorial_2a( valor ):
    # fatorial de 0 é 1 por definição
    if valor == 0:
        return 1
    fatorial = 1
    for numero in range( 1, valor+1 ):
        # ou podemos fazer fatorial*=numero
        fatorial = fatorial*numero
    return fatorial
    
def calcular_power_2b( valor, potencia ):
    n = 1
    saida = 1
    while n <= potencia:
        saida *= valor
        n+=1
    return saida

def calcular_exponencial_2c( x, m_termos ):
    n = 0
    saida = 0
    while n <= m_termos:
        xn = calcular_power_2b(x , n)
        n_fatorial = calcular_fatorial_2a( n )
        saida += xn/n_fatorial
        n+=1
    return saida
\end{minted}
}

    \vspace{0.5cm}
    \item A sequência de Fibonacci é uma série numérica em que cada termo é dado pela soma dos dois anteriores, iniciando em 0 e 1:
    \[
    0, 1, 1, 2, 3, 5, 8, 13, 21, 34, 55, \dots
    \]

    \noindent Escreva uma função em Python que receba como argumento um número inteiro $n$ ($n > 0$), exiba na tela a sequência até o $n$-ésimo termo e retorne o valor do $n$-ésimo elemento. A função deve validar o parâmetro $n$ antes da execução do cálculo, retornando os seguintes códigos de erro:
    \begin{itemize}
    \item \code{-1} caso $n$ \textbf{não seja} do tipo inteiro (\code{int});
    \item \code{-2} caso $n$ \textbf{não seja} estritamente maior que zero ($n \le 0$).
    \end{itemize}

    \noindent  Utilize a estrutura de repetição \code{for} para fazer a questão.

    \ifmostrarGabarito{
\begin{minted}{python}
# Resposta:
def calcular_fibonacci( n ):

    if not type(n)==int:
        return -1

    if n <= 0:
        return -2

    soma = 0
    a = 0
    b = 1
    for termo in range( 1 , n+1 , 1 ):
        
        if termo == 1:
            valor = 0
        elif termo == 2:
            valor = 1
        else:
            valor = a+b
            a = b
            b = valor
        print(f"O termo {n} da sequencia é {valor}")
            
        soma+=valor
    return valor
\end{minted}
}

    \vspace{0.5cm}
    \item Escreva uma função em Python que receba dois números inteiros positivos, $n$ e $s$. A função deve calcular a soma dos números pares no intervalo de $0$ até $n$, interrompendo o processo imediatamente antes que o somatório ultrapasse o limite $s$.
    A função deve validar os parâmetros $n$ e $s$ antes do cálculo, retornando os seguintes valores:
    \begin{itemize}
        \item \code{-1} caso $n$ ou $s$ \textbf{não sejam} do tipo inteiro (\code{int});
        \item \code{-2} caso $n$ ou $s$ \textbf{não sejam} estritamente positivos ($n \le 0$ ou $s \le 0$);
        \item O valor do somatório resultante, caso os dados de entrada sejam válidos.
    \end{itemize}

    \noindent É \textbf{obrigatório} utilizar as estruturas \code{for}, \code{continue} e \code{break} na solução.

    \ifmostrarGabarito{
\begin{minted}{python}
# Resposta:
def somar_pares( n, s ):

    if not (type(n)==int) or not (type(s)==int):
        print("Os valores de n e s precisam ser inteiros.")
        return -2
        
    if n < 0 or s < 0:
        print("Os valores de n e s precisam ser positivos.")
        return -2

    soma = 0
    for numero in range( 0, n+1, 1 ):

        if n % 2 != 0:
            continue
            
        # devo testar antes de contabilizar
        if soma+numero > s:
            break

        # como não passou o limite s, eu contabilizo
        soma += numero
    return soma        

\end{minted}
}
\end{enumerate}